% TOP_PROBLEM: {"problemId": "problem.existence-of-one-way-permutations", "canonicalTitle": "Existence of One-Way Permutations", "releaseRank": 50, "primaryDomain": "cryptography_coding_information_optimization", "primaryDomainLabel": "Cryptography, coding, information & optimization", "displayStatus": "open", "releaseStatus": "open"}
% !TeX program = lualatex
% Definition and sources only; this file is not an LLM solution attempt.
\documentclass[11pt]{article}
\usepackage[margin=25mm]{geometry}
\usepackage{fontspec}
\setmainfont{Latin Modern Roman}
\usepackage{amsmath,amssymb,mathtools}
\usepackage{unicode-math}
\setmathfont{Latin Modern Math}
\usepackage{xurl,hyperref}
\hypersetup{colorlinks=true,urlcolor=blue}
\setcounter{secnumdepth}{0}
\setlength{\parindent}{0pt}
\setlength{\parskip}{5pt}
\setlength{\emergencystretch}{3em}
\begin{document}
\section*{\texorpdfstring{Existence of One-Way Permutations}{Existence of One-Way Permutations}}\label{50}
\subsection{Definitions and mathematical statement}
For every $n\ge1$, require a bijection $f_n:\{0,1\}^n\to\{0,1\}^n$, with one polynomial-time algorithm evaluating $f_n(x)$ from $(1^n,x)$. For each probabilistic polynomial-time $A$, require
\[
 \Pr_{x\leftarrow\{0,1\}^n,A}[A(1^n,f_n(x))=x]\le\eta_A(n),
\]
where $\eta_A$ is negligible. Since $f_n$ is bijective, each output has exactly one preimage. This is an unconditional existence problem on all binary strings of each length, without sampled keys or variable domains.
\par\smallskip\textit{Formulation and concept references:} \href{https://eprint.iacr.org/2015/752.pdf}{[S1]}.

\subsection{Short English statement}
Can a length-preserving bijection be easy to compute but negligibly likely to be inverted by any efficient randomized algorithm?
\subsection{Sources}
\begin{itemize}
\item[S1]On Constructing One-Way Permutations from Indistinguishability Obfuscation, ePrint2015/752, TCC2016 version.
\url{https://eprint.iacr.org/2015/752.pdf}.
\textit{Formulation source.}
\item[S2]On One-Shot Signatures, Quantum vs Classical Binding, and Obfuscating Permutations.
\url{https://arxiv.org/abs/2507.12456}.
\textit{Further reference; not a proof endorsement.}
\item[S3]Quantum Search with In-Place Queries.
\url{https://drops.dagstuhl.de/storage/00lipics/lipics-vol350-tqc2025/LIPIcs.TQC.2025.1/LIPIcs.TQC.2025.1.pdf}.
\textit{Further reference; not a proof endorsement.}
\end{itemize}
\end{document}
